% page 439 of the facsimile (image p-438 of the scan)
% block 170.2 x 242.0 mm ; left margin 31.8 mm
\pgc{10.160mm}{0.000mm}{
\l{\cel{52}431}
\l{}
\l{\sou{pendulo}.\cel{2}Stango tenua, o filo, finanta per maso, qua su movas}
\l{\cel{3}ye plano vertikala, e di qua la ocili esas izokrona. - DEFIRS.}
\l{}
\l{\sou{penetrar}.\cel{2}I. (\sou{netrans.}) Enirar transirante to quo esas obstaklo. -}
\l{\cel{3}II. (\sou{trans.}) Enirar kelke profunde aden (ulo). - DEFIS.}
\l{}
\l{\sou{peninsulo}.\cel{2}(\sou{geogr.}) Granda lando, cirkondata da aquo an omna}
\l{\cel{3}lateri minus una per qua lu juntesas a la kontinento. - DEFIS.}
\l{}
\l{\sou{peniso}. (\sou{anat.})\cel{2}Organo cilindroida,vaskuloza, rigideskiva, uzata}
\l{\cel{3}por la kopulaco e koito che la mamiferi. - DEFIS.}
\l{}
\l{\sou{penitencar}. (\sou{netrans.}) (religio). I. Repentar pri pekir. - II.}
\l{\cel{3}Expiacar pro pekir. - EFIS.}
\l{}
\l{\sou{pensar}. (\sou{trans. e netrans.})\cel{2}(\sou{ulo, ad ulo, pri ulu od ulo)}.-}
\l{\cel{3}I. Uzar sua spirito por konceptar, judikar,ulo. II. Uzar}
\l{\cel{3}sua fakultato konceptar, judikar. Posedar ta fakultato. - EFIS.}
\l{}
\l{\sou{pensionar}, (\sou{trans.})\cel{2}Pagar sen-intermanke, sive a stat-employato,}
\l{\cel{3}sive por lojeyo, sive por nutreso, sive por edukeso od instrukt-}
\l{\cel{3}eso. - DEFIRS.}
\l{}
\l{\sou{pento}. Direciono di plano, qua formacas angulo obliqua ye hori-}
\l{\cel{3}zontalo. - FIS.}
\l{}
\l{\sou{penta}- . Sufixo qua signifikas \textquotedbl{}qua havas kin...\textquotedbl{}.}
\l{}
\l{\sou{pentaedro.}\cel{2}(\sou{geom.}) Solido qua havas kin edri. - DEFIS.}
\l{}
\l{\sou{pentagono}. (geom.) Figuro geometriala qua havas kin anguli e kin}
\l{\cel{3}lateri. - DEFIS.}
\l{}
\l{\sou{pentagramo}. Figuro talismana qua reprezentas stelo kin-pinta.}
\l{\cel{3}- DEF.}
\l{}
\l{\sou{pentametr}o. Verso de kin pedi. - DEFIRS.}
\l{}
\l{\sou{pentatlo}. (en\cel{1}\sou{la epoki antiqua)}.\cel{2}Asemblajo de la kin exerci}
\l{\cel{3}atletala (lukto, kuro, salto, lanso di disko, lanso di jave-}
\l{\cel{3}lino). - DEFIS.}
\l{}
\l{\sou{pentekosto}. (\sou{religio krist.)}\cel{2}I. (\sou{che la Judi)}\cel{1}. Festo memorige}
\l{\cel{3}di la lego donita sur la monto Sinai, celebrata sep semani pos}
\l{\cel{3}la Pasko-dio duesma, e dum qua onu fris a Deo la premici di la}
\l{\cel{3}rekolto. - II. (\sou{che la}\cel{1}kristani). Festo celebrata en la sundio}
\l{\cel{3}sepesma pos Pasko, memorige di la decenso di la Santa Spirito}
\l{\cel{3}sur la apostoli. - EFIS.}
\l{}
\l{\sou{pecnio}. (\sou{bot.}) Planto ranunkulacea, kun granda flori intense reda,}
\l{\cel{3}bunta o blanka. - L.\cel{1}\sou{paeonia}. - DEFIS.}
\l{}
\l{pepito.\cel{2}(mineral.) Mikra maso de metalo quan onu renkontras nature}
\l{\cel{3}sen gango. - DFIS.}
}
